\documentclass[11pt]{article}
\usepackage[utf8]{inputenc}
\usepackage[T1]{fontenc}
\usepackage{amsmath,amssymb,amsthm}
\usepackage[margin=1in]{geometry}
\newcommand{\leg}[2]{\left(\frac{#1}{#2}\right)}
\DeclareMathOperator{\Imag}{Im}
\DeclareMathOperator{\Real}{Re}
\newtheorem{theorem}{Theorem}
\newtheorem{lemma}{Lemma}
\title{The closed form for $\displaystyle S_n=\sum_{\ell=0}^{n-1}(\ell^2\bmod n)$\\[2pt]
\large a step-by-step derivation}
\author{}
\date{}
\begin{document}
\maketitle

\section{The quantity and the goal}

For an integer $n\ge2$ we study
\[
S_n=\sum_{\ell=0}^{n-1}\bigl(\ell^2\bmod n\bigr),
\]
where $\ell^2\bmod n\in\{0,1,\dots,n-1\}$ is the least nonnegative remainder of $\ell^2$ on division by $n$.

\subsection{An equivalent ``by value'' description}
Instead of summing over $\ell$, we may sum over the possible \emph{values} $j=\ell^2\bmod n$. If
$N(j)=\#\{x\bmod n:\ x^2\equiv j\}$ is the number of square roots of $j$, then collecting equal values,
\[
S_n=\sum_{j=0}^{n-1} j\,N(j).
\]
This viewpoint explains where the arithmetic enters: $N(j)$ is a multiplicative, character-theoretic object,
and $S_n$ is its first moment. We will not use this form directly, but it is the reason class numbers appear.

\subsection{Notation fixed once and for all}
Throughout, set
\[
\zeta=e^{2\pi i/n},\qquad n=m^2 d\ \ (d\text{ squarefree}),\qquad
G(t)=\sum_{\ell=0}^{n-1}\zeta^{t\ell^2}.
\]
Here $m^2$ is the largest square dividing $n$, so $m=\prod_{p^a\,\|\,n}p^{\lfloor a/2\rfloor}$; and $G(t)$ is a
\emph{quadratic Gauss sum}. We will prove:

\begin{theorem}\label{thm:main}
\[
\boxed{\;S_n=\frac{n(n-m)}{2}
-2n\sum_{\substack{\Delta<0\ \text{a discriminant}\\ \Delta\,\mid\, n}}\frac{h(\Delta)}{w(\Delta)}\;}
\]
where $\Delta$ runs over all \emph{negative discriminants} dividing $n$ --- that is, all
$\Delta\equiv0,1\pmod4$, $\Delta<0$, with $-\Delta\mid n$ (fundamental and non-fundamental alike) ---
$h(\Delta)$ is the class number of the order $\mathcal{O}_\Delta$ of discriminant $\Delta$ (equivalently the
number of $\mathrm{SL}_2(\mathbb{Z})$-classes of primitive positive-definite binary quadratic forms of
discriminant $\Delta$), and $w(\Delta)=|\mathcal{O}_\Delta^\times|\in\{2,4,6\}$ is its number of units.
\end{theorem}

\noindent The proof first produces the equivalent expanded form over \emph{fundamental} discriminants only,
\begin{equation}\label{eq:double}
S_n=\frac{n(n-1)}{2}-\frac{n(m-1)}{2}
-2n\sum_{\substack{D<0\ \text{fundamental}\\ D\,\mid\, n}}\frac{h_{\mathbb{Q}(\sqrt D)}}{w_D}
\sum_{D\cdot s^2\,\mid\, n}s\prod_{p\mid s}\Bigl(1-\tfrac1p\leg{D}{p}\Bigr),
\end{equation}
in which $D$ runs over negative fundamental discriminants dividing $n$, $h_{\mathbb{Q}(\sqrt D)}$ is the class
number and $w_D=|\mathcal{O}_{\mathbb{Q}(\sqrt D)}^\times|$ the number of units of $\mathbb{Q}(\sqrt D)$, and
$\leg{D}{p}$ is the Kronecker symbol; \S\ref{sec:collapse} then collapses the double sum into the single sum of
Theorem~\ref{thm:main}.

\section{Strategy}

The obstacle is the operation ``$\bmod n$,'' which is not algebraic. The plan is to \emph{linearise} it:
\begin{enumerate}
\item Expand the residue function $a\mapsto(a\bmod n)$ in the additive characters $\zeta^{ta}$ of $\mathbb{Z}/n$
      (finite Fourier analysis). This converts the awkward ``$\bmod$'' into a clean sum of exponentials.
\item Substituting $a=\ell^2$ and summing over $\ell$ replaces $\sum_\ell\zeta^{t\ell^2}$ by the Gauss sum $G(t)$,
      a thoroughly understood object.
\item Gauss's evaluation of $G(t)$ produces a \emph{quadratic character}, and Dirichlet's class number formula
      turns the remaining sum into class numbers.
\end{enumerate}
Steps 1--2 are elementary; the two deep classical theorems live in Step~3.

\section{The Fourier identity}

\subsection{Finite Fourier analysis on $\mathbb{Z}/n$}
Every function $f:\mathbb{Z}/n\to\mathbb{C}$ has a unique expansion
\[
f(a)=\sum_{t=0}^{n-1}\hat f(t)\,\zeta^{ta},\qquad \hat f(t)=\frac1n\sum_{a=0}^{n-1}f(a)\,\zeta^{-ta},
\]
because the characters $\{\zeta^{ta}\}_{t}$ are orthogonal:
$\sum_{a=0}^{n-1}\zeta^{(t-t')a}=n$ if $t\equiv t'$ and $0$ otherwise (a geometric series whose ratio
$\zeta^{t-t'}$ is an $n$-th root of unity $\ne1$ when $t\not\equiv t'$). We apply this to the specific function
$f(a)=a\bmod n$, i.e.\ $f(a)=a$ for $a\in\{0,\dots,n-1\}$.

\subsection{A geometric-sum lemma}
We need the coefficients $\hat f(t)$, which require $\sum_{a=0}^{n-1}a\,\omega^a$ for a root of unity $\omega$.

\begin{lemma}\label{lem:geo}
If $\omega^n=1$ and $\omega\ne1$, then $\displaystyle\sum_{a=0}^{n-1}a\,\omega^{a}=\frac{n}{\omega-1}$.
\end{lemma}

\begin{proof}
Differentiate the geometric sum $\sum_{a=0}^{n-1}x^a=\dfrac{x^n-1}{x-1}$ and multiply by $x$:
\[
\sum_{a=0}^{n-1}a\,x^{a}=\frac{n x^{n}(x-1)-x(x^{n}-1)}{(x-1)^2}.
\]
Set $x=\omega$ and use $\omega^n=1$: the numerator becomes $n(\omega-1)-\omega\cdot0=n(\omega-1)$, so the right
side is $n(\omega-1)/(\omega-1)^2=n/(\omega-1)$.
\end{proof}

\subsection{Fourier coefficients of the residue function}
The $t=0$ coefficient is the average value, $\hat f(0)=\frac1n\sum_{a=0}^{n-1}a=\frac{n-1}{2}$. For $t\ne0$,
Lemma~\ref{lem:geo} with $\omega=\zeta^{-t}$ gives
\[
\hat f(t)=\frac1n\sum_{a=0}^{n-1}a\,\zeta^{-ta}=\frac1n\cdot\frac{n}{\zeta^{-t}-1}=\frac{1}{\zeta^{-t}-1}.
\]
Therefore, for every $a$,
\begin{equation}\label{eq:resid}
a\bmod n=\frac{n-1}{2}+\sum_{t=1}^{n-1}\frac{\zeta^{ta}}{\zeta^{-t}-1}.
\end{equation}
This is the key linearisation: the left side is a ``sawtooth,'' the right side a finite trigonometric sum.

\subsection{Summing over $\ell$: enter the Gauss sum}
Put $a=\ell^2$ in \eqref{eq:resid} and sum over $\ell=0,\dots,n-1$. The constant term contributes
$n\cdot\frac{n-1}2=\binom n2$, and interchanging the two finite sums,
\begin{equation}\label{eq:S-gauss}
S_n=\binom n2+\sum_{t=1}^{n-1}\frac{1}{\zeta^{-t}-1}\sum_{\ell=0}^{n-1}\zeta^{t\ell^2}
=\binom n2+\sum_{t=1}^{n-1}\frac{G(t)}{\zeta^{-t}-1}.
\end{equation}
The inner sum is precisely the Gauss sum $G(t)$. Everything arithmetic is now hidden in $G(t)$.

\subsection{The cotangent identity}
We simplify the prefactor $\dfrac{1}{\zeta^{-t}-1}$. Writing $\theta=\pi t/n$ so that $\zeta^{-t}=e^{-2i\theta}$,
and multiplying numerator and denominator by $e^{i\theta}$,
\[
\frac{1}{e^{-2i\theta}-1}=\frac{e^{i\theta}}{e^{-i\theta}-e^{i\theta}}=\frac{\cos\theta+i\sin\theta}{-2i\sin\theta}
=-\frac12+\frac{i}{2}\cot\theta.
\]
Hence
\begin{equation}\label{eq:cot}
\frac{1}{\zeta^{-t}-1}=-\frac12+\frac{i}{2}\cot\frac{\pi t}{n}.
\end{equation}
The real constant $-\tfrac12$ will produce an elementary term; the imaginary $\cot$ part will produce the
class-number term.

\subsection{Separating the two contributions}
Insert \eqref{eq:cot} into \eqref{eq:S-gauss}:
\begin{equation}\label{eq:two-sums}
S_n-\binom n2=-\frac12\sum_{t=1}^{n-1}G(t)\;+\;\frac{i}{2}\sum_{t=1}^{n-1}G(t)\cot\frac{\pi t}{n}.
\end{equation}
We treat the two sums separately.

\subsubsection*{The constant sum}
Summing $G(t)$ over \emph{all} $t$ and swapping order,
\[
\sum_{t=0}^{n-1}G(t)=\sum_{\ell=0}^{n-1}\sum_{t=0}^{n-1}\zeta^{t\ell^2}
=\sum_{\ell=0}^{n-1} n\,[\,n\mid \ell^2\,]=n\,N_0,\qquad N_0:=\#\{x\bmod n:\ x^2\equiv0\},
\]
using $\sum_{t=0}^{n-1}\zeta^{tk}=n$ when $n\mid k$ and $0$ otherwise. Since $G(0)=n$, we get
$\sum_{t=1}^{n-1}G(t)=nN_0-n$, so the first sum in \eqref{eq:two-sums} equals
\begin{equation}\label{eq:A}
A:=-\tfrac12\sum_{t=1}^{n-1}G(t)=\frac{n(1-N_0)}{2}.
\end{equation}

\subsubsection*{The cotangent sum: only the imaginary part survives}
Two symmetries collapse the second sum. First, replacing $\ell$ by anything does not matter, but replacing
$t$ by $n-t$ does:
\[
G(n-t)=\sum_{\ell}\zeta^{(n-t)\ell^2}=\sum_\ell\zeta^{-t\ell^2}=\overline{G(t)},\qquad
\cot\frac{\pi(n-t)}{n}=\cot\!\Bigl(\pi-\frac{\pi t}{n}\Bigr)=-\cot\frac{\pi t}{n}.
\]
Now split $G(t)=\Real G(t)+i\,\Imag G(t)$. Under $t\mapsto n-t$, $\Real G$ is \emph{even} and $\Imag G$ is
\emph{odd}, while $\cot$ is odd. Thus $\Real G(t)\cot\frac{\pi t}{n}$ is odd in $t$ and sums to $0$, whereas
$\Imag G(t)\cot\frac{\pi t}{n}$ is even and survives:
\[
\sum_{t=1}^{n-1}G(t)\cot\frac{\pi t}{n}=i\sum_{t=1}^{n-1}\Imag G(t)\,\cot\frac{\pi t}{n}.
\]
Multiplying by $\tfrac i2$ turns the leading $i$ into $-\tfrac12$, so the second sum in \eqref{eq:two-sums} is
\begin{equation}\label{eq:B}
B:=-\frac12\sum_{t=1}^{n-1}\Imag G(t)\,\cot\frac{\pi t}{n}.
\end{equation}

\subsection{The master decomposition}
Combining \eqref{eq:two-sums}, \eqref{eq:A}, \eqref{eq:B}:
\begin{equation}\label{eq:master}
\boxed{\;S_n-\binom n2=A+B,\qquad A=\frac{n(1-N_0)}{2},\quad B=-\frac12\sum_{t=1}^{n-1}\Imag G(t)\,\cot\frac{\pi t}{n}.\;}
\end{equation}
Everything elementary is in $A$; everything deep is in $B$.

\section{The term $A$: the squareful correction}

We must count $N_0=\#\{x\bmod n:\ x^2\equiv0\}$. By the Chinese Remainder Theorem this factors over the prime
powers $p^a\,\|\,n$, and $x^2\equiv0\pmod{p^a}$ exactly when $p^{\lceil a/2\rceil}\mid x$. The number of such $x$
modulo $p^a$ is $p^{a}/p^{\lceil a/2\rceil}=p^{\lfloor a/2\rfloor}$. Multiplying,
\[
N_0=\prod_{p^a\,\|\,n}p^{\lfloor a/2\rfloor}=m.
\]
Hence
\begin{equation}\label{eq:Aval}
A=\frac{n(1-m)}{2}=-\frac{n(m-1)}{2}.
\end{equation}
In words: the squares put $m$-fold extra weight on the residue $0$, and this is the only purely elementary
deviation from the average. It vanishes precisely when $n$ is squarefree ($m=1$).

\section{The term $B$: the class-number correction}

This is the heart of the proof. We evaluate
$\sum_{t=1}^{n-1}\Imag G(t)\cot\frac{\pi t}{n}$ in five steps.

\subsection{Step 1 --- reduction to primitive Gauss sums}
Let $g=\gcd(t,n)$ and write $n=gM$, $t=gc$ with $\gcd(c,M)=1$. In $G(t)=\sum_{k=0}^{n-1}e^{2\pi i gck^2/(gM)}
=\sum_{k=0}^{n-1}e^{2\pi i ck^2/M}$ the summand depends only on $k\bmod M$, and as $k$ ranges over
$0,\dots,gM-1$ each residue modulo $M$ occurs $g$ times. Therefore
\[
G(t)=g\,g(c;M),\qquad g(c;M):=\sum_{k=0}^{M-1}e^{2\pi i ck^2/M},\qquad
\cot\frac{\pi t}{n}=\cot\frac{\pi c}{M}.
\]
So every $G(t)$ is a rational multiple of a \emph{primitive} Gauss sum $g(c;M)$, $\gcd(c,M)=1$, where
$M=n/\gcd(t,n)$ is a divisor of $n$. This is the natural ``denominator'' of the angle $\pi t/n$.

\subsection{Step 2 --- Gauss's sign theorem}
The exact value of $g(c;M)$ --- including the notoriously subtle \emph{sign} --- was determined by Gauss. For
$\gcd(c,M)=1$,
\[
g(c;M)=\begin{cases}
\sqrt M\,\leg{c}{M} & M\equiv1\pmod4,\\[1mm]
i\,\sqrt M\,\leg{c}{M} & M\equiv3\pmod4,\\[1mm]
0 & M\equiv2\pmod4,\\[1mm]
(1+i^{\,c})\,\leg{M}{c}\,\sqrt M & M\equiv0\pmod4.
\end{cases}
\]
Taking imaginary parts,
\[
\Imag g(c;M)=\begin{cases}0 & M\equiv1,2\pmod4,\\[1mm]
\sqrt M\,\leg{c}{M} & M\equiv3\pmod4,\\[1mm]
\sqrt M\,\leg{-M}{c} & M\equiv0\pmod4.\end{cases}
\]
\textbf{The crucial observation:} in every nonzero case $\Imag g(c;M)=\sqrt M\,\chi_D(c)$, where
$\chi_D=\leg{D}{\cdot}$ is an \emph{odd, real, primitive} Dirichlet character --- the Kronecker character of an
imaginary quadratic field $\mathbb{Q}(\sqrt D)$, $D<0$ a fundamental discriminant. (Concretely $D=-M$ when $M$ is
squarefree $\equiv3\pmod4$; for $4\mid M$ the symbol $\leg{-M}{\cdot}$ reduces to the character of the
squarefree-kernel discriminant.) This is where the field $\mathbb{Q}(\sqrt D)$ is born: it is forced on us by the
sign of the Gauss sum.

\subsection{Step 3 --- which modulus $M$ feeds which discriminant}
A character $\chi_D$ has \emph{conductor} $|D|$ (its smallest period). Comparing conductors on both sides of
$\Imag g(c;M)=\sqrt M\,\chi_D(c)$ shows that the modulus $M$ produces $\chi_D$ exactly when $M$ is the conductor
times a perfect square:
\[
M=|D|\,s^2,\qquad s\ge1.
\]
Intuitively: the ``squarefull part'' $s^2$ of $M$ is invisible to the primitive character $\chi_D$ (it only
affects which residues are coprime to $M$), so it can be any square, while the squarefree skeleton of $M$ is
pinned to the conductor $|D|$.

\subsection{Step 4 --- grouping the sum}
We now organise $\sum_{t}\Imag G(t)\cot\frac{\pi t}{n}$ by the modulus $M=n/\gcd(t,n)$, then by $D$. For a fixed
$M\mid n$, the $t$'s with $n/\gcd(t,n)=M$ are exactly $t=\tfrac nM c$ with $\gcd(c,M)=1$; for these
$G(t)=\tfrac nM\,g(c;M)$ and $\cot\frac{\pi t}{n}=\cot\frac{\pi c}{M}$. Hence the $M$-block contributes
$\frac nM\sum_{\gcd(c,M)=1}\Imag g(c;M)\cot\frac{\pi c}{M}$. Substituting Step~2 and Step~3,
\[
\sum_{t=1}^{n-1}\Imag G(t)\cot\frac{\pi t}{n}
=\sum_{\substack{D<0\ \text{fund.}\\ D\mid n}}\ \sum_{|D|s^2\mid n}
\frac{n}{|D|s^2}\,\sqrt{|D|s^2}\;T_D(s),
\]
where we abbreviate the \emph{finite cotangent sum at level $M=|D|s^2$}
\[
T_D(s):=\sum_{\substack{c\bmod |D|s^2\\ \gcd(c,\,|D|s^2)=1}}\chi_D(c)\,\cot\frac{\pi c}{|D|s^2}.
\]

\subsection{Step 5 --- Dirichlet's class number formula}
Two classical facts evaluate $T_D(s)$.

\paragraph{The value of $L(1,\chi_D)$.}
Dirichlet's class number formula states that, for the character $\chi_D$ of the imaginary quadratic field
$\mathbb{Q}(\sqrt D)$,
\[
L(1,\chi_D)=\frac{2\pi\,h_{\mathbb{Q}(\sqrt D)}}{w_D\,\sqrt{|D|}}\;>\;0,
\qquad L(s,\chi_D)=\sum_{k\ge1}\frac{\chi_D(k)}{k^s}.
\]
Here $h_{\mathbb{Q}(\sqrt D)}$ is the class number and $w_D$ the number of units. The positivity
$L(1,\chi_D)>0$ is the single genuinely deep input of the whole proof; no elementary proof of it is known.

\paragraph{The finite cotangent sum as an $L$-value.}
For an odd character, $L(1,\cdot)$ has a closed finite expression through cotangents. The level-$M$ sum
$T_D(s)$ realises $L(1,\psi_M)$ for the imprimitive character $\psi_M$ modulo $M=|D|s^2$ induced from $\chi_D$;
passing from $\psi_M$ to the primitive $\chi_D$ removes the Euler factors at the primes of $M$ outside the
conductor:
\[
T_D(s)=\frac{2M}{\pi}\,L(1,\psi_M)
=\frac{2M}{\pi}\,L(1,\chi_D)\!\!\prod_{\substack{p\mid M\\ p\nmid|D|}}\!\!\Bigl(1-\frac{\chi_D(p)}{p}\Bigr).
\]
Substituting $M=|D|s^2$ and the value of $L(1,\chi_D)$,
\[
T_D(s)=\frac{2|D|s^2}{\pi}\cdot\frac{2\pi\,h_{\mathbb{Q}(\sqrt D)}}{w_D\sqrt{|D|}}\!\!
\prod_{\substack{p\mid s\\ p\nmid|D|}}\!\!\Bigl(1-\frac{\chi_D(p)}{p}\Bigr)
=\frac{4\sqrt{|D|}\,s^2\,h_{\mathbb{Q}(\sqrt D)}}{w_D}\!\!\prod_{\substack{p\mid s\\ p\nmid|D|}}\!\!\Bigl(1-\frac{\chi_D(p)}{p}\Bigr).
\]
(The primes $p\mid M$ outside the conductor are exactly the primes $p\mid s$ with $p\nmid|D|$.) Finally, since
$\chi_D(p)=\leg{D}{p}=0$ whenever $p\mid D$, the factor for such $p$ is $1$, so the restriction $p\nmid|D|$ is
vacuous and the product may run over \emph{all} $p\mid s$.

\subsection{Assembly of $B$}
Put $\sqrt{|D|s^2}=\sqrt{|D|}\,s$ and the formula for $T_D(s)$ into the grouped sum of Step~4. The $(D,s)$ term is
\[
\frac{n}{|D|s^2}\cdot\sqrt{|D|}\,s\cdot\frac{4\sqrt{|D|}\,s^2\,h_{\mathbb{Q}(\sqrt D)}}{w_D}\prod_{p\mid s}\Bigl(1-\tfrac1p\leg Dp\Bigr)
=\frac{4n\,h_{\mathbb{Q}(\sqrt D)}}{w_D}\,s\prod_{p\mid s}\Bigl(1-\tfrac1p\leg Dp\Bigr),
\]
where the two powers of $\sqrt{|D|}$ and of $s$ cancel against the denominator $|D|s^2$. Therefore
\[
\sum_{t=1}^{n-1}\Imag G(t)\cot\frac{\pi t}{n}
=\sum_{\substack{D<0\ \text{fund.}\\ D\mid n}}\frac{4n\,h_{\mathbb{Q}(\sqrt D)}}{w_D}\sum_{D\cdot s^2\mid n}s\prod_{p\mid s}\Bigl(1-\tfrac1p\leg Dp\Bigr),
\]
and, using $B=-\tfrac12(\cdots)$ from \eqref{eq:B},
\begin{equation}\label{eq:Bval}
B=-2n\sum_{\substack{D<0\ \text{fund.}\\ D\mid n}}\frac{h_{\mathbb{Q}(\sqrt D)}}{w_D}\sum_{D\cdot s^2\mid n}s\prod_{p\mid s}\Bigl(1-\tfrac1p\leg Dp\Bigr).
\end{equation}

\section{Collapsing the double sum: from fields to orders}\label{sec:collapse}

The expanded form \eqref{eq:double} sums over fundamental discriminants $D$ and, for each, over conductors $s$
with $Ds^2\mid n$. We now recognise the inner summand as a single class number and merge the two sums into one.

\subsection{The inner summand is a class number of an order}
Fix a fundamental discriminant $D<0$ and a conductor $s\ge1$, and let $\mathcal{O}_{Ds^2}\subset\mathcal{O}_D$ be the
order of conductor $s$ inside the maximal order $\mathcal{O}_D$ of $\mathbb{Q}(\sqrt D)$; it has discriminant
$\Delta=Ds^2$. The class-number formula for non-maximal orders (Cox, \emph{Primes of the form $x^2+ny^2$},
Thm.~7.24) reads
\[
h(Ds^2)=\frac{h_{\mathbb{Q}(\sqrt D)}\;s}{[\mathcal{O}_D^\times:\mathcal{O}_{Ds^2}^\times]}
\prod_{p\mid s}\Bigl(1-\tfrac1p\leg Dp\Bigr).
\]
The unit index is exactly the ratio of unit counts, $[\mathcal{O}_D^\times:\mathcal{O}_{Ds^2}^\times]=w_D/w(Ds^2)$
(here $w(Ds^2)=2$ for $s>1$, while $w_D\in\{2,4,6\}$). Substituting and dividing by $w(Ds^2)$,
\begin{equation}\label{eq:order-id}
\frac{h_{\mathbb{Q}(\sqrt D)}}{w_D}\;s\prod_{p\mid s}\Bigl(1-\tfrac1p\leg Dp\Bigr)
=\frac{h(Ds^2)}{w(Ds^2)}.
\end{equation}
Thus the entire $(D,s)$ summand of \eqref{eq:double} is $h(\Delta)/w(\Delta)$ for the discriminant
$\Delta=Ds^2$; the conductor product and the exceptional units $w_D\in\{4,6\}$ are absorbed into the single
object $h(\Delta)/w(\Delta)$.

\subsection{Re-indexing by discriminant}
Every negative discriminant $\Delta$ factors \emph{uniquely} as $\Delta=Ds^2$ with $D$ fundamental and $s\ge1$
its conductor. Hence the pairs $(D,s)$ with $D$ fundamental, $D\mid n$ and $Ds^2\mid n$ are in bijection with the
negative discriminants $\Delta\mid n$ (note $Ds^2\mid n\iff\Delta\mid n$, which forces $D\mid n$ automatically).
Applying \eqref{eq:order-id} term by term to \eqref{eq:Bval},
\begin{equation}\label{eq:Bsingle}
B=-2n\sum_{\substack{D<0\ \text{fund.}\\ D\mid n}}\frac{h_{\mathbb{Q}(\sqrt D)}}{w_D}\sum_{D\cdot s^2\mid n}s\prod_{p\mid s}\Bigl(1-\tfrac1p\leg Dp\Bigr)
=-2n\sum_{\substack{\Delta<0\ \text{disc.}\\ \Delta\mid n}}\frac{h(\Delta)}{w(\Delta)}.
\end{equation}

\section{Conclusion}

Inserting $A$ from \eqref{eq:Aval} and $B$ from \eqref{eq:Bsingle} into the master decomposition
\eqref{eq:master}, and condensing the elementary part via
$\binom n2-\frac{n(m-1)}2=\frac{n(n-1)}2-\frac{n(m-1)}2=\frac{n(n-m)}2$, gives
\[
S_n=\frac{n(n-m)}{2}
-2n\sum_{\substack{\Delta<0\ \text{disc.}\\ \Delta\mid n}}\frac{h(\Delta)}{w(\Delta)},
\]
which is Theorem~\ref{thm:main}. The intermediate equality \eqref{eq:double} over fundamental discriminants is
\eqref{eq:Bsingle} read right-to-left. \qed

\subsection*{Two remarks}
\emph{An inequality for free.} Every class number satisfies $h(\Delta)\ge1$, every unit count $w(\Delta)>0$, and
$m\ge1$; hence both corrections are $\ge0$ and $S_n\le\binom n2$, with equality iff $n$ is squarefree and has no
divisor $\equiv3\pmod4$ (equivalently, no negative discriminant divides $n$). The same positivity is visible in
the expanded form \eqref{eq:double}, where each Euler factor
$1-\tfrac1p\leg Dp\in\{1-\tfrac1p,\,1,\,1+\tfrac1p\}$ is positive.

\emph{Where the depth sits.} Steps 1--4 of $B$, and all of $A$, are elementary; the collapse of
\S\ref{sec:collapse} is pure bookkeeping via the order class-number formula. Only two classical theorems are
used as black boxes: Gauss's sign evaluation of $g(c;M)$ (\S5.2) and Dirichlet's class number formula
(\S5.5). The lone non-elementary fact inside them is $L(1,\chi_D)>0$.
\end{document}
